\documentclass{amsart}
\usepackage{amssymb}
\renewcommand{\baselinestretch}{1.5}
\addtolength{\textwidth}{.2in}
\addtolength{\topmargin}{-.5in}
\addtolength{\textheight}{1in}

\usepackage{ifthen}
\usepackage{graphicx}
\usepackage{xcolor}


\newlength{\cwidth}
\newcommand{\cents}{\settowidth{\cwidth}{c}%
\divide\cwidth by2
\advance\cwidth by-.1pt
c\kern-\cwidth
\vrule width .1pt depth.2ex height1.2ex
\kern\cwidth}

\newcommand{\sageprompt}{ {\tt sage$>$} }
\newcommand{\tab}{\rule{20pt}{0pt}}
\newcommand{\blnk}{\rule{1.5pt}{0pt}\rule{.4pt}{1.2pt}\rule{9pt}{.4pt}\rule{.4pt}{1.2pt}\rule{1.5pt}{0pt}}
\newcommand{\suchthat}{\; \rule[-3pt]{.25pt}{13pt} \;}
\newcommand{\divides}{\!\mid\!}
\newcommand{\tdiv}{\; \mbox{div} \;}
\newcommand{\restrict}[2]{#1 \,\rule[-4pt]{.125pt}{14pt}_{\,#2}}
\newcommand{\lcm}[2]{\mbox{lcm} (#1, #2)}
\renewcommand{\gcd}[2]{\mbox{gcd} (#1, #2)}
\newcommand{\Naturals}{{\mathbb N}}
\newcommand{\Integers}{{\mathbb Z}}
\newcommand{\Znoneg}{{\mathbb Z}^{\mbox{\tiny noneg}}}
\newcommand{\Enoneg}{{\mathbb E}^{\mbox{\tiny noneg}}}
\newcommand{\Qnoneg}{{\mathbb Q}^{\mbox{\tiny noneg}}}
\newcommand{\Rnoneg}{{\mathbb R}^{\mbox{\tiny noneg}}}
\newcommand{\Rationals}{{\mathbb Q}}
\newcommand{\Reals}{{\mathbb R}}
\newcommand{\Complexes}{{\mathbb C}}
%\newcommand{\F2}{{\mathbb F}_{2}}
\newcommand{\relQ}{\mbox{\textsf Q}}
\newcommand{\relR}{\mbox{\textsf R}}
\newcommand{\nrelR}{\mbox{\raisebox{1pt}{$\not$}\rule{1pt}{0pt}{\textsf R}}}
\newcommand{\relS}{\mbox{\textsf S}}
\newcommand{\relA}{\mbox{\textsf A}}
\newcommand{\Dom}[1]{\mbox{Dom}(#1)}
\newcommand{\Cod}[1]{\mbox{Cod}(#1)}
\newcommand{\Rng}[1]{\mbox{Rng}(#1)}

\DeclareMathOperator\caret{\raisebox{1ex}{$\scriptstyle\wedge$}}

\newtheorem*{defi}{Definition}
\newtheorem*{exer}{Exercise}
\newtheorem{thm}{Theorem}[section]
\newtheorem*{thm*}{Theorem}
\newtheorem{lem}[thm]{Lemma}
\newtheorem{cor}{Corollary}
\newtheorem{conj}{Conjecture}

\renewenvironment{proof}%
{\begin{quote} \emph{Proof:} }%
{\rule{0pt}{0pt} \newline \rule{0pt}{15pt} \hfill Q.E.D. \end{quote}}


\addtolength{\abovedisplayskip}{0pt}
\addtolength{\belowdisplayskip}{24pt}
\addtolength{\abovedisplayshortskip}{0pt}
\addtolength{\belowdisplayshortskip}{48pt}

\newcommand{\hs}{\rule{18pt}{0pt}}
\newcommand{\vs}{\rule[-16pt]{0pt}{32pt}}

\begin{document}
\thispagestyle{empty}

\centerline{\Large Activity 26 -- Introduction to Proof}
\centerline{\large formulas for sums and products}

\bigskip
\Large


\begin{enumerate}
\item Complete the following tables about sums of squares.
\bigskip

\begin{tabular}{c|c}
\vs\hs $n$ \hs & \hs $\displaystyle \sum_{j=1}^n \, j^2$ \hs \\ \hline
\vs $1$ & $1$ \\
\vs $2$ & \\
\vs $3$ & \\
\vs $4$ & \\
\end{tabular}
\hspace{.5in}
\begin{tabular}{c|c}
\vs \hs $n$ \hs & \hs $\displaystyle \frac{n(n + 1)(2n + 1)}{6}$ \hs \\ \hline
\vs $1$ & $1$ \\
\vs $2$ & \\
\vs $3$ & \\
\vs $4$ & \\
\end{tabular}

\vspace{.3in}

\item The base case for an inductive proof that

\[ \sum_{j=1}^n \, j^2 \; = \; \frac{n(n + 1)(2n + 1)}{6}\]

really just consists of noticing that the first rows in the two tables
above are identical. Write the base case formally.

\vfill

\newpage

\item The inductive step in a proof that

\[ \sum_{j=1}^n \, j^2 \; = \; \frac{n(n + 1)(2n + 1)}{6}\]

involves verifying the following equality

\[ (k + 1)^2 \, + \, \frac{n(n + 1)(2n + 1)}{6} \; = \; 
\frac{(k + 1)((k + 1) + 1)(2(k + 1) + 1)}{6}
\]

or (after simplifying the right-hand side)

\[ (k + 1)^2 \, + \, \frac{n(n + 1)(2n + 1)}{6} \; = \; 
\frac{(k + 1)(k + 2)(2k + 3)}{6}
\]

\noindent Do it! (Try to take advantage of the factored form as much as
possible.)

\vfill

\newpage

\item Put the last problems together and write a complete inductive proof
of the formula for the sum of squares.

\vfill

\newpage

\item The image below illustrates that a perfect square can be decom-
posed into a sum of odd numbers. (The L-shaped things are a
visual representation for odd numbers. Do you see why?)
\bigskip

\centerline{ \begin{picture}(0,0)%
\includegraphics{./L-decomp.pdf}%
\end{picture}%
\setlength{\unitlength}{3947sp}%
%
\begingroup\makeatletter\ifx\SetFigFont\undefined%
\gdef\SetFigFont#1#2#3#4#5{%
  \reset@font\fontsize{#1}{#2pt}%
  \fontfamily{#3}\fontseries{#4}\fontshape{#5}%
  \selectfont}%
\fi\endgroup%
\begin{picture}(2049,1824)(1189,-2773)
\end{picture}%
 }
\bigskip

Notice that the number of terms in the sum is the quantity that is
squared, so

\[ \forall n \in \Naturals, \; \sum_{i=0}^n \, 2i+1 \; = \; (n+1)^2 \]

\noindent (Since a sum that goes from $0$ to $n$ has $n + 1$ terms.)

\noindent What will we need to add to both sides of the inductive hypothesis?

\vspace{.75in}

Work out the algebraic part in the inductive step:

\[(n + 1)^2 + \underline{\hspace{.75in}} \; = \; ((n + 1) + 1)^2 \] 

\vfill

\newpage

\item The triangular numbers are the numbers that can be expressed as
triangular arrays of points.
\bigskip

\centerline{ \begin{picture}(0,0)%
\includegraphics{./triangular_numbers.pdf}%
\end{picture}%
\setlength{\unitlength}{3947sp}%
\begin{picture}(5963,1088)(1193,-2116)
\end{picture}%
 }
\bigskip

Let $T_n$ represent the $n$-th triangular number. So $T_1 = 1$, $T_2 = 3$,
$T_3 = 6$, {\em et cetera}.

Each triangular number can be thought of as a sum, which allows
us to deduce a formula for $T_n$.  What is it?

\vfill

\item Find a formula for the sum

\[ \sum_{j=1}^n \; T_j, \]

\noindent and prove it using mathematical induction.

\vfill

\newpage

\item A {\em hexagonal number} is one that can be represented by a hexagonal
array of dots. A formula for the hexagonal numbers is

\[ H_n \, = \, 3n^2 - 3n + 1. \]

Find a formula for the sum

\[
\sum_{j=1}^n \, H_j 
\] 

\noindent and prove it using mathematical induction.

\vfill

\newpage

\item (Exercise 5 in GIAM \S 5.2) Prove the following formula for a prod-
uct.

\[ \prod_{i=2}^n \, \left( 1 - \frac{1}{i} \right) \; = \; \frac{1}{n}  \]

\vfill

\newpage

\item Here is a formula for the sum of the first $n$ fourth powers

\[ \frac{n \cdot (n + 1) \cdot (2n + 1) \cdot (3n^2 + 3n - 1)}{30}. \]

Prove it using induction.

\vfill


\end{enumerate}


\end{document}
